\subsection{Is Affect Necessary for Moral Concern?}
\label{sec:2.3.4}
The claim that reason does not manufacture concern from nothing is an empirical claim, and the evidence usually produced for it is thinner than the use it gets put to. Section~\ref{sec:5.3.1} sets out the standard exhibit: the footbridge case engages emotional processing that the switch case leaves quiet, and patients with ventromedial prefrontal damage endorse the push at a higher rate than people with intact emotional response \autocite{greene2001fmri}, \autocite{koenigs2007damage}. That establishes that affect changes particular moral judgments. It does not establish that affect is what makes anything matter, which is the stronger claim this book needs, and the gap between the two is where an argument has to be made rather than gestured at.

Five lines of evidence bear on the stronger claim, and one serious argument runs the other way.

\runin{Affect is part of valuation, not a coloring applied to it}
The clearest case is not about morality at all. Patients with ventromedial prefrontal damage retain intelligence, language, working memory, and explicit reasoning about consequences, and they choose disastrously: on a card task with hidden payoff structures, they keep returning to decks that pay well immediately and punish severely over time, long after control subjects have learned to avoid them, and long after some of them can state the rule out loud \autocite{bechara1994insensitivity}. Knowing which deck is bad and being steered by that knowledge come apart. Whatever the affective signal is doing in a healthy chooser, it is not decorating a decision that reasoning had already made; the reasoning survives its loss and the decision does not.

\runin{The philosophical position is empirical and has been argued as such}
Moral sentimentalism — the claim that moral judgment is constitutively tied to affective response — is usually met as Hume's, and its modern form is a psychological hypothesis with evidence attached. Shaun Nichols argues for what he calls sentimental rules: normative rules and affective mechanisms make separate contributions to moral judgment, and the rules that carry the distinctive force of moral prohibition are the ones bound to an emotional response to the harm they prohibit \autocite{nichols2004sentimental}. The interesting part for a designer is the separability. Rules without the affective component do not stop being rules; they stop being the kind of thing that carries the force.

\runin{The prosocial emotions are motivational states, not readouts}
Compassion has been characterized as a distinct emotion with a specific functional profile: it arises at the perception of another's suffering and moves the agent toward the sufferer, rather than reporting the suffering to some further faculty that decides what to do \autocite{goetz2010compassion}. Section~\ref{sec:2.2.1}'s training study is the mechanistic version of the same finding, with compassion practice recruiting reward and approach circuitry where empathy practice recruited shared distress \autocite{klimecki2013differential}. Guilt has the same shape pointed at one's own conduct: guilt, which construes the act as wrong, predicts reparative behavior, while shame, which construes the self as wrong, predicts withdrawal and defensiveness \autocite{tangney2007moral}. In each case the emotion is doing motivational work directly. It is not information that motivation then acts on.

\runin{Intact reasoning coexists with impaired concern, and the literature disagrees about how}
The dissociation this book needs is the one where moral knowledge survives and moral concern does not, and it is the most contested item on the list. James Blair's finding was that children and adults with psychopathic traits fail to distinguish moral transgressions from conventional ones — treating hitting another child and talking out of turn as wrong in the same way, wrong because prohibited — which he attributed to a deficit in responding to distress cues \autocite{blair1995cognitive}. Later work reversed the surface finding: on a battery of moral dilemmas, incarcerated psychopaths gave the same judgments as controls, which the authors summarized in their title as knowing right from wrong and not caring \autocite{cima2010psychopaths}.

Those two results are not easily reconciled, and I am not going to pretend the field has done it. What survives both is the shape that matters here. On Blair's account the affective deficit degrades the moral judgment itself; on the later account the judgment is intact and the connection to conduct is severed. Either way the case is a dissociation between competence and agency in exactly section~\ref{sec:2.3}'s terms, and neither account contains an instance of concern generated by reasoning in the absence of affect.

\runin{The rationalist objection, which is the strongest thing said against all of this}
Recognizing a reason might itself motivate. On a Kantian account, moral agency runs on the recognition of what one has reason to do, and affect is neither the source of that recognition nor required for it to move anyone. Jeanette Kennett makes the version of this argument that engages the evidence rather than bypassing it: high-functioning autistic people, whose empathic responses are atypical or absent, are nonetheless conscientious moral agents, and the capacity that carries their moral seriousness is a reverence for reason rather than a felt resonance with others \autocite{kennett2002autism}. If that is right, the empirical generalization above has a standing counterexample, and it is a counterexample in the only species available for study.

I take the argument seriously and do not think it closes the gap, for one reason. Reverence for reason is itself a caring — about consistency, about getting it right, about not being the sort of agent who makes exceptions for himself. Kennett's autistic agents are described as moved by moral considerations, and being moved is the explanandum. The case shows convincingly that \emph{empathy} is not the required mechanism, which section~\ref{sec:2.2.3} had already concluded for independent reasons. It does not show that a system in which nothing at all registers as mattering can be moved by a normative relation it has computed.

\runin{What this licenses and what it does not}
Perhaps non-affective moral motivation is possible. Producing a verbal justification or computing a normative relation is not evidence that the consideration matters to the system, and that is the epistemic problem entire: the output of a system that cares and the output of a system that has learned what caring sounds like are the same text. Every agent known to care is an affective agent, and no design has yet demonstrated concern arising from reasoning alone.

That is not a proof that non-affective moral agency is metaphysically impossible, and section~\ref{sec:3.4} does not treat it as one. It places the evidential burden on whoever proposes non-affective concern as an engineering plan, and a burden of that kind is what an engineering decision is normally made under. Section~\ref{sec:3.3} weighs how much that burden actually carries, and reduces it: the machine half of the record reports on a route nobody has attempted, which is weaker evidence than an absence usually is.
